

\section{Partitionnement par intervalle}

\begin{frame}
\frametitle{Principe}

La collection est \alert{triée} sur la clé  et on la découpe en fragments d'une taille maximale pré-déterminée (de 64MO à qq GO)

\vfill

 Chaque fragment couvre donc un intervalle $[min_f, max_f[$.
 
\vfill
 
\figSlideWithSize{partition-intervalle}{10cm}
\end{frame}


\begin{frame}
\frametitle{Dynamicité: comment évolue le système\,?}

L'opération de base est le \alert{split}, division d'un gros fragment en deux petits.

\figSlideWithSize{split}{11cm}

\alert{L'équilibrage} consiste à répartir équitablement les fragments.

\end{frame}



\begin{frame}
\frametitle{Efficacité du routage}


% A two-columns presentation
\begin{tabular}{l|p{6cm}}
   \parbox{5cm}{
   
   Pour 1 TO de données, et une taille d'entrée  $(I, a)$ de 20 octets.

\begin{redbullet}
  \item Fragment de 4 KO:
    le routage décrit 250 millions de fragments, soit une taille de 5 GO;    
  \item Fragment de 1 MO: il faudra 1 million de fragments,
    et seulement 20 MO pour la structure de routage.
\end{redbullet}
Dans tous les  cas, le routage tient en mémoire RAM

   }  % Inclusion of an XML file
  & 
  \parbox{6cm}{
 
\begin{tabular}{r|l}
Intervalle  & Serveur\\
\hline
$]-\infty, a]$ & A  \\
$[a+1, b]$ & B  \\
$[b+1, c]$ & C  \\
$[c+1, d]$ & B  \\
$[d+1, \infty[$ & A  \\
\hline
\end{tabular}

}

\end{tabular}



 

\end{frame}


\begin{frame}
\frametitle{Les opérations du routeur}

% A two-columns presentation
\begin{tabular}{l|p{6cm}}
   \parbox{5cm}{
$get(k)$: chercher l'intervalle qui contient $k$, transmettre
au serveur correspondant.

\smallskip

$put(k, d)$: identifier le serveur avec $k$, transmettre l'insertion
 
 \smallskip
 
 $delete(k)$: idem
 
 \smallskip
 
 $range(k_1, k_2)$: transmettre à tous les serveurs gérant un intervalle chevauchant $[k_1, k_2]$
 
 \smallskip
 
 Toute autre opération\,: transmettre à tous les serveurs.
   }  % Inclusion of an XML file
  & 
  \parbox{6cm}{
 
\begin{tabular}{r|l}
Intervalle  & Serveur\\
\hline
$]-\infty, a]$ & A  \\
$[a+1, b]$ & B  \\
$[b+1, c]$ & C  \\
$[c+1, d]$ & B  \\
$[d+1, \infty[$ & A  \\
\hline
\end{tabular}

}

\end{tabular}

\end{frame}


\section{Etude de cas\,: \mongodb/}


\begin{frame}
\frametitle{Le \textit{sharding} de MongoDB}

Le routeur connaît les intervalles, et les serveurs associés.

\figSlideWithSize{mongo-shard}{10cm}

\vfill

\begin{redbullet}
   \item Chaque fragment est un \textit{replica set} complet.
   \item On peut ajouter ou supprimer un serveur dynamiquement.
\end{redbullet}

\end{frame}


\begin{frame}
\frametitle{Partitionnement et équilibrage}

\figSlideWithSize{partition-split}{11cm}

\end{frame}



\begin{frame}[fragile]
\frametitle{Allons-y pour un petit test}

On lance un \textit{config server}

\begin{minted}{text}
mkdir /data/configdb
mongod --configsvr --dbpath /data/configdb --port 27019
\end{minted}

On lance un routeur en lui indiquant où se trouve(nt) le(s) \textit{config server}(s)

\begin{minted}{text}
mongos --configdb localhost:27019
\end{minted}

Et deux serveurs (il faudrait des \textit{replica sets} complets en production)
\begin{minted}{text}
mongod --dbpath /data/node1 --port=30001 --oplogSize 700
mongod --dbpath /data/node2 --port=30002 --oplogSize 700
\end{minted}

À partir de \textit{mongos}, déclarons les deux serveurs comme \textit{shards}.

\begin{minted}{text}
mongo --port 27017
sh.addShard("localhost:30000")
sh.addShard("localhost:30001")
sh.enableSharding("nfe204")
\end{minted}

\end{frame}

\begin{frame}[fragile]
\frametitle{Choisir la clé de partitionnement}

On indique la clé de partitionnement avec la commande suivante\,:

\begin{minted}{text}
sh.shardCollection("nfe204.videos", { "title": 1 } )
\end{minted}

Attention au choix de la clé

\begin{redbullet}
  \item Les clés-séquences\,:  on envoie toujours les insertions au même serveur\,; bon ou pas?
  \item Les clés à faible cardinalité (ex: pays)\,: pas très bon non plus, car 
               il pourra être difficile de faire un \textit{split}
\end{redbullet}

Bonne option: appliquer un hachage pour bien distribuer les insertions\,:

\begin{minted}{text}
sh.shardCollection("nfe204.videos", { "title": "hashed" } )
\end{minted}

\end{frame}


\begin{frame}[fragile]
\frametitle{Inspecter la configuration}

Liste des \textit{shards}
\begin{minted}{text}
  mongos> db.runCommand({listshards: 1})
\end{minted}

Statut général du cluster
\begin{minted}{text}
    mongos> sh.status()    
    mongos> db.stats()
    mongos> db.printShardingStatus()
\end{minted}

Les fragments sont dans une collection \textit{chuncks}.

\begin{minted}{text}
use config
db.chuncks.count()
db.chuncks.findOne()
\end{minted}
\end{frame}



\begin{frame}[fragile]
\frametitle{Chargement en masse}

Utiliser un générateur de données:
ipsum (\url{https://github.com/10gen-labs/ipsum}).

\vfill

Charger 1 million de pseudo-films.

\begin{minted}{text}
python ./pymonipsum.py -d nfe204 
           -c movies ---count 1000000 movies.jsch
\end{minted}

\vfill

Surveiller le chargement et les éclatements...

\end{frame}
